\input{preamble}

\begin{document}

\pagestyle{fancy}
\fancyhead[L]{Seconde}
\fancyhead[C]{\textbf{Fiche préparation DS n°6}}
\fancyhead[R]{\today}

Faire cette fiche \textbf{sans calculatrice}. L'utiliser uniquement pour vérifier vos résultats !

\exe{}{
Sans utiliser de calculatrice, déterminer la parité des nombres ci-dessous :
\begin{multicols}{2}
\begin{enumerate}[label=(\alph*)]
\item $129~767 + 60~012$
\item $340~841 - 122~195$
\item $2^5 \times 3^8 \times 7^9 \times 11^4$
\item $27 \times 42~657 + 48 \times 51~444$
\end{enumerate}
\end{multicols}
}{exe:1}{
\begin{multicols}{2}
\begin{enumerate}[label=(\alph*)]
\item impair
\item pair
\item pair
\item impair
\end{enumerate}
\end{multicols}
}

\exe{}{
	Décomposer les nombres suivants en produits de facteurs premiers.
	\begin{multicols}{2}
	\begin{enumerate}[label=\roman*)]
		\item  $360$
		\item $165$
		\item $12~100$
		\item  $144 \times 9$
	\end{enumerate}
	\end{multicols}
}{exe:2}{

	\begin{multicols}{2}
	\begin{enumerate}[label=\roman*)]
		\item $360 = 2^3 \times 3^2 \times 5$
		\item $165 = 3 \times 5 \times 11$
		\item $12~100 = 2^2 \times 5^2 \times 11^2$
		\item $144 \times 9 = 2^4 \times 3^4$
	\end{enumerate}
	\end{multicols}
}

\exe{}{
	Exprimer chaque valeur suivante sous forme de fraction irréductible.
	\begin{multicols}{2}
	\begin{enumerate}[label=\alph*), itemsep=10pt]
		\item $\dfrac{105}{63}$
		\item $\dfrac{3~150}{5~250}$
		\item $1 - \dfrac{24}{42}$
		\item $\dfrac{44}{147} \times \dfrac{63}{55}$
	\end{enumerate}
	\end{multicols}
}{exe:3}{
	\begin{multicols}{2}
	\begin{enumerate}[label=\alph*), itemsep=10pt]
		\item $\dfrac{105}{63} = \dfrac53$
		\item  $\dfrac{3~150}{5~250} = \dfrac35$
		\item $\dfrac{3}{7}$
		\item $\dfrac{44}{147} \times \dfrac{63}{55} = \dfrac{12}{35}$
	\end{enumerate}
	\end{multicols}
}

\exe{}{
Justifier que les nombres suivants sont premiers ou pas.
    \begin{multicols}{2}
\begin{enumerate}[itemsep=2em]
	\item \begin{minipage}[t]{\linewidth} 327 \end{minipage}
	\item \begin{minipage}[t]{\linewidth} 130 \end{minipage}
	\item \begin{minipage}[t]{\linewidth} 155 \end{minipage}
	\item \begin{minipage}[t]{\linewidth} 97 \end{minipage}
\end{enumerate}
\end{multicols}
}{exe:4}{
\begin{enumerate}[itemsep=1em]
	\item Comme 3 + 2 + 7 = 12 est un multiple de 3 donc 327 aussi, il admet donc au moins trois diviseurs qui sont 1, 3 et lui-même,
	\\{\bfseries \color[HTML]{f15929}327 n'est donc pas premier}.
	\item Comme 130 est pair, il admet donc au moins trois diviseurs qui sont 1, 2 et lui-même, 
	{\bfseries \color[HTML]{f15929}130 n'est donc pas premier}.
	\item Comme le chiffre des unités de 155 est un 5, alors 155 est divisible par 5, il admet donc au moins trois diviseurs qui sont 
	1, 5 et lui-même, {\bfseries \color[HTML]{f15929}155 n'est donc pas premier}.
	\item En effectuant la division euclidienne de 97 par tous les nombres premiers inférieurs à $\sqrt{97}$, c'est-à-dire par les
	 nombres 2, 3, 5 et 7, le reste n'est jamais nul, {\bfseries \color[HTML]{f15929}97 est donc un nombre premier}.
\end{enumerate}
}

\exe{}{
	Donner $\D_{48}$ et en déduire les entiers naturels $x\in\N$ vérifiant $x^2 - 8x= 48$.
	
	\emph{Indication : montrer d'abord que $x^2 - 8x = x(x-8)$.}
}{exe:5}{
	$\D_{48} = \bigset{1 ; 48 ; 2 ; 24 ; 3 ; 16 ; 4 ; 12 ; 6 ; 8}$.

	Soit $x\in\N$ vérifiant $x^2 - 8x = 48 \iff x(x-8) = 48$.
	Comme $x\in\N$ est entier, $x-8$ l'est aussi, et les deux nombres divisent 48.
	
	On cherche donc deux éléments de $\D_{48}$ dont le produit vaut 48 et qui sont séparés de $8$ unités.
	Par inspection, seul $4\times12 = 48$ et donc $x=12$ fonctionne, et c'est la seule solution dans $\N$ de l'équation.
	
	Par symétrie, il existe une autre solution, $x = -4$, qui appartient à $\Z$ mais pas à $\N$.
}

%%%%%%%%%%%%

\newpage
\fancyhead[C]{\textbf{Solutions}}
\shipoutAnswer

\end{document}
